% GERADO por tools/generation/gerar_secoes_latex.py a partir de
% artifacts/paper/secoes/04-metodo.md. Nao editar a mao: a proxima geracao
% desfaz. Corrija o Markdown e regere.
O envelope metodológico é o Design Science Research, que trata a pesquisa como construção e avaliação de artefatos para resolver problemas de classe identificada \cite{hevner2004design,peffers2007design}. A escolha não abre um ciclo novo: o próprio MPO é artefato de um ciclo de DSR em curso, construído a partir de estudos exploratórios \cite{vieira2020universal}, publicado em versão preliminar \cite{vieira2021model} e evoluído por avaliações com especialistas e estudos de caso publicados nos Anais do SBSI \cite{vieira2022evaluating} e por survey \cite{vieira2023survey}, até sua versão consolidada \cite{vieira2022thesis}. Este trabalho é nova iteração desse ciclo: submete o modelo a um critério que as avaliações anteriores não aplicaram --- a consistência de seus compromissos ontológicos --- e devolve ao ciclo um modelo revisado com a evidência que justifica cada revisão.

O método de construção do artefato ontológico é a Methontology \cite{fernandez1997methontology}, que organiza o desenvolvimento em especificação, conceitualização, formalização, integração, implementação e avaliação, com artefatos intermediários prescritos para cada etapa. Duas alternativas foram consideradas: a NeOn \cite{suarez2012neon} organiza o desenvolvimento em cenários centrados no reuso de recursos ontológicos e em redes de ontologias, e sua força não é exercida aqui, onde a fonte é um único documento em prosa; a LOT \cite{poveda2022lot} elicita requisitos junto a interlocutores em ciclos iterativos, e pressupõe disponibilidade de especialistas, que está fora do escopo deste trabalho (\S{}8). A Methontology foi mantida por duas razões. A primeira é a rastreabilidade: seus artefatos intermediários --- glossário de termos, árvore de classificação de conceitos, dicionário de verbos e tabelas de fórmulas e regras --- ligam cada termo do documento em prosa ao elemento formal que dele derivou, e é essa cadeia que torna auditável o traço de cada deficiência até o ponto do modelo que a permitiu. A segunda é a comparabilidade: a formalização analisada foi construída sob a mesma metodologia, e assim o antes e o depois se comparam etapa a etapa.

O procedimento da análise teve quatro passos. \textbf{Reconstrução:} a formalização publicada foi remontada em OntoUML como linha de base, a partir dos diagramas em resolução de origem e dos artefatos de conceitualização, sem nenhuma correção --- corrigi-la destruiria seu valor como termo de comparação. \textbf{Leitura:} cada classe, estereótipo e relação da linha de base foi confrontado com o trecho do modelo de referência que o originou. \textbf{Qualificação:} contou como deficiência todo ponto que satisfez dois testes simultâneos --- o modelo formal se compromete com algo que a UFO proíbe ou deixa de distinguir, \textit{e} esse compromisso não é imposto pelo texto do modelo de referência. Sem o primeiro teste o achado seria preferência de modelagem; sem o segundo, conserto de erro próprio. \textbf{Classificação:} cada achado recebeu uma das quatro categorias da Representation Theory \cite{wand1993ontological} quando é deficiência de mapeamento, ou foi marcado como fora da tipologia, com a razão registrada. As correções foram aplicadas em duas rodadas --- a primeira estrutural, a segunda de adoção das microteorias de eventos e de agentes ---, preservadas em separado, de modo que linha de base e rodadas sejam comparáveis entre si. O modelo revisado foi transformado em OWL pela transformação gUFO oficial \cite{almeida2019gufo}, com customização aditiva sobre o gerado e publicação do \textit{diff} (\S{}6).

A avaliação foi desenhada em três frentes, com pesos distintos. A primeira, e única que valida as afirmações substantivas do artigo, é a \textbf{instanciação seguida de consulta}: a ontologia revisada foi povoada com os dados de um observatório de projetos de pesquisa e extensão de uma universidade pública brasileira, documentado nos Anais Estendidos do SBSI \cite{sbsi_estendido}, e interrogada por sete questões de competência em SPARQL. As sete são questões de domínio --- quem executou uma carga e quando, qual a cadeia de proveniência de um conteúdo divulgado, se dois observatórios que se dizem aderentes ao modelo cobrem os mesmos conceitos ---, e não perguntas sobre a arquitetura da ontologia, que qualquer ontologia responderia do mesmo modo; a comparação entre iniciativas exigiu um segundo observatório, sintético, declarado como limitação (\S{}8). A segunda frente é a \textbf{verificação de conformidade à UFO} sobre a linha de base e o modelo revisado, que atesta conformidade sintática e semântica, não substância; a terceira é a \textbf{detecção de \textit{pitfalls} na OWL} \cite{poveda2014oops}, igualmente executada antes e depois. As duas últimas só produzem evidência como diferença entre execuções, e por isso cada verificador foi submetido a controle positivo --- execução deliberada sobre um modelo com defeito conhecido, que prova que a regra dispara ---, sem o que um relatório vazio não distinguiria \textit{sem problemas} de \textit{regra que nunca dispara}.
